\begin{abstract}
We validate that layer-wise weight tokenization preserves sufficient structural signal for architecture family classification, achieving 80\% $\pm$ 4.2\% test accuracy on 100 pretrained vision models (ResNet, ViT, EfficientNet, ConvNeXt)---significantly exceeding our 60\% validation gate and 25\% random baseline. The tokenization approach---flattening weight matrices, zero-padding to fixed length, and applying per-layer normalization---enables both a simple baseline (per-layer statistics: mean/std/norm) and a transformer encoder (cross-layer attention) to reach this performance. Surprisingly, per-layer statistics match transformer accuracy despite 10$\times$ fewer parameters (200K vs 2M), revealing that architecture families exhibit strong layer-level separability on small datasets (70 training samples). Transformer overfitting (100\% train, 73.33\% $\pm$ 4.2\% val) indicates capacity calibration is critical in weight-space learning. Building on prior weight-space property prediction work, our controlled comparison of tokenization strategies establishes layer-wise processing as a validated foundation for weight-space transformers while highlighting that simple statistical baselines provide competitive performance on coarse-grained tasks. Future work should explore larger datasets and harder tasks (property prediction, backdoor detection) where cross-layer reasoning may provide measurable benefits.

\textbf{Keywords}: Weight-space learning, neural network tokenization, meta-learning, architecture classification, transformers
\end{abstract}
